\titulo{2. Marco teórico profundizado}

\subtitulo{2.1. Tamaño de entrada $n$}
El análisis de un algoritmo comienza definiendo qué magnitud representa el tamaño de la entrada, porque la función de costo solo tiene sentido cuando $n$ está asociado de forma consistente con aquello que crece en el problema. En esta práctica, $n$ representa el número de registros de \texttt{CustomerID} almacenados en la muestra, incluidos identificadores repetidos; no representa la cantidad de clientes distintos. Esta elección permite comparar la búsqueda lineal y la búsqueda binaria sobre colecciones del mismo tamaño y evita interpretar de forma ambigua una expresión como $T(n)$. Cormen et al. presentan el tamaño de entrada como la variable básica respecto de la cual se expresa el tiempo de ejecución, mientras Levitin y Skiena resaltan que su definición depende del problema concreto y de la representación elegida (Cormen et al., 2022; Levitin, 2012; Skiena, 2020).

La unidad utilizada para $n$ también determina qué conclusiones son válidas. Por ejemplo, si se midiera número de clientes únicos en lugar de registros, las muestras con identificadores repetidos tendrían un tamaño diferente aun cuando el algoritmo recorriera exactamente la misma cantidad de posiciones de memoria. Por ello, en análisis de algoritmos conviene separar la dimensión lógica del problema de otras características de los datos, como distribución, duplicados o valor numérico de las claves. Esta disciplina permite comparar implementaciones sin modificar silenciosamente la variable independiente del experimento (Cormen et al., 2022; Sedgewick \& Wayne, 2011; McGeoch, 2012).

\subtitulo{2.2. Modelo RAM}
El modelo RAM (Random Access Machine) abstrae una computadora ideal en la que operaciones elementales como acceso a memoria, asignación, comparación y aritmética básica se contabilizan con costo constante. Su utilidad no consiste en reproducir exactamente lo que ocurre en un procesador moderno, sino en proporcionar un lenguaje común para razonar sobre la cantidad de trabajo que realiza un algoritmo al crecer la entrada. Gracias a esta abstracción se puede comparar la estructura del costo de dos soluciones sin depender de la frecuencia de CPU, el compilador o el sistema operativo de una máquina particular (Cormen et al., 2022; Levitin, 2012; Knuth, 1997).

El modelo RAM exige declarar un convenio de conteo. Una misma línea de código puede descomponerse en accesos, comparaciones, sumas y asignaciones, y dos análisis pueden producir constantes diferentes sin contradecirse si ambos mantienen un criterio coherente. Lo relevante para el análisis asintótico es que las diferencias de costo constante no alteran el orden de crecimiento dominante. En esta práctica se conserva además un contador de sondeos separado del tiempo real, de manera que la evidencia estructural del algoritmo no se confunda con la velocidad de una plataforma específica (Cormen et al., 2022; Skiena, 2020; Zanabria Gálvez, 2026).

\subtitulo{2.3. Función de costo $T(n)$}
La función $T(n)$ representa el número de operaciones elementales ejecutadas por un algoritmo bajo un modelo de costo previamente definido. No es, por sí misma, una medición en segundos: es una expresión matemática que permite observar cómo cambia el trabajo al aumentar $n$. Derivar $T(n)$ obliga a identificar qué instrucciones se ejecutan una vez, cuáles dependen del número de iteraciones y cuáles solo aparecen bajo determinadas condiciones. Esta derivación es un puente entre el código concreto y las cotas asintóticas que luego resumen su comportamiento (Cormen et al., 2022; Levitin, 2012; Skiena, 2020).

Para búsqueda lineal, por ejemplo, la cantidad de posiciones inspeccionadas puede variar desde una hasta $n$. En un peor caso ausente, el ciclo debe visitar cada posición; si se adopta un convenio RAM que contabiliza inicializaciones, control del ciclo, acceso, comparación, actualización de contadores e incremento del índice, se obtiene una expresión de la forma $an+b$. El valor concreto de $a$ y $b$ depende del convenio, pero la dependencia lineal con $n$ se mantiene. En búsqueda binaria, en cambio, la cantidad de intervalos posibles se reduce aproximadamente a la mitad en cada iteración, lo que conduce a una función relacionada con $\log_2 n$ (Cormen et al., 2022; Sedgewick \& Wayne, 2011; Levitin, 2012).

\subtitulo{2.4. Operación dominante y tasas de crecimiento}
Cuando $n$ aumenta, algunos términos de $T(n)$ influyen mucho más que otros. Si $T(n)=7n+5$, el término $7n$ domina el crecimiento y la constante aditiva deja de ser significativa para entradas grandes. De manera similar, en una expresión como $3n^2+4n+20$, el término cuadrático domina a los lineales y constantes. Esta simplificación no afirma que las constantes sean irrelevantes en tiempos reales; afirma que, desde la perspectiva asintótica, el tipo de crecimiento está determinado por el término que aumenta más rápidamente (Cormen et al., 2022; Levitin, 2012; Knuth, 1997).

Las tasas habituales $1$, $\log n$, $n$, $n\log n$ y $n^2$ permiten ordenar algoritmos según su escalabilidad teórica. Para tamaños pequeños, una implementación con peor orden asintótico puede ser más rápida por constantes favorables, menor sobrecarga o mejor localidad de memoria. Sin embargo, al crecer suficientemente la entrada, las diferencias entre tasas de crecimiento terminan dominando. Por eso el análisis asintótico se complementa, pero no se reemplaza, con mediciones empíricas (Cormen et al., 2022; Skiena, 2020; McGeoch, 2012).

\subtitulo{2.5. Notación $O$}
La notación $O(g(n))$ expresa una cota superior asintótica. Formalmente, $T(n)\in O(g(n))$ si existen constantes positivas $c$ y $n_0$ tales que $0\le T(n)\le c\,g(n)$ para todo $n\ge n_0$. La notación permite afirmar que, a partir de un tamaño suficientemente grande, el crecimiento de $T(n)$ no supera un múltiplo constante de $g(n)$. Por ejemplo, $7n+5\in O(n)$ porque puede acotarse por un múltiplo de $n$ para todo $n$ suficientemente grande (Cormen et al., 2022; Levitin, 2012; Skiena, 2020).

Un error frecuente consiste en usar $O$ como sinónimo automático de ``peor caso''. $O$ describe una relación entre funciones y puede utilizarse para acotar una función de mejor caso, promedio o peor caso, siempre que se especifique cuál se está analizando. Por ejemplo, el mejor caso de una búsqueda lineal que encuentra la clave en la primera posición tiene costo constante y pertenece a $O(1)$, mientras que su peor caso pertenece a $O(n)$. La distinción entre cota y escenario es central en la guía de esta práctica (Cormen et al., 2022; Black, 2020; Zanabria Gálvez, 2026).

\subtitulo{2.6. Notación $\Omega$}
La notación $\Omega(g(n))$ expresa una cota inferior asintótica. Se cumple cuando existen constantes positivas $c$ y $n_0$ tales que $T(n)\ge c\,g(n)$ para todo $n\ge n_0$. Mientras $O$ establece un techo asintótico, $\Omega$ establece un piso para la función analizada. Esta herramienta es útil para demostrar que una cantidad de trabajo no puede crecer por debajo de cierto orden dentro del modelo adoptado (Cormen et al., 2022; Levitin, 2012; Skiena, 2020).

También aquí es incorrecto identificar $\Omega$ con ``mejor caso'' de manera automática. El peor caso de búsqueda lineal, cuando la clave está ausente, no solo es $O(n)$: también es $\Omega(n)$ porque debe inspeccionar proporcionalmente $n$ posiciones. Por tanto, ese mismo peor caso es $\Theta(n)$. La notación inferior describe la función de costo elegida, no el tipo de entrada que produjo dicha función (Cormen et al., 2022; Black, 2020; Zanabria Gálvez, 2026).

\subtitulo{2.7. Notación $\Theta$ y cota ajustada}
La notación $\Theta(g(n))$ expresa una cota ajustada: la función está acotada superior e inferiormente por múltiplos positivos de $g(n)$ a partir de cierto $n_0$. En términos formales, existen $c_1,c_2>0$ y $n_0$ tales que $c_1g(n)\le T(n)\le c_2g(n)$ para todo $n\ge n_0$. Esto permite describir el orden de crecimiento con mayor precisión que una cota superior aislada (Cormen et al., 2022; Levitin, 2012; Skiena, 2020).

Decir que el peor caso de búsqueda binaria es $\Theta(\log n)$ significa que su número de iteraciones crece del mismo orden que un logaritmo, tanto por arriba como por abajo, bajo el modelo y escenario definidos. No significa que $\Theta$ sea sinónimo de ``caso promedio''. Un promedio puede ser $\Theta(n)$, $\Theta(\log n)$ o cualquier otro orden dependiendo del algoritmo y de la distribución de entradas. Separar estas dos ideas evita una de las confusiones conceptuales más comunes al comenzar el análisis asintótico (Cormen et al., 2022; Sedgewick \& Wayne, 2011; Zanabria Gálvez, 2026).

\subtitulo{2.8. Mejor, promedio y peor caso}
Los casos describen conjuntos o distribuciones de entradas, no tipos de notación. En búsqueda lineal, una clave ubicada en la primera posición produce el mejor escenario porque se necesita un solo sondeo; una clave ubicada al final o ausente produce el peor escenario porque se inspeccionan $n$ posiciones. El caso promedio requiere declarar supuestos probabilísticos, por ejemplo que todas las posiciones de éxito sean equiprobables, en cuyo caso el número esperado de sondeos es $(n+1)/2$ (Cormen et al., 2022; Levitin, 2012; Skiena, 2020).

En experimentación es importante no confundir la repetición de una clave con un promedio teórico. Ejecutar 30 veces la misma consulta permite observar variabilidad temporal bajo condiciones controladas, pero no transforma esa consulta en una muestra de todas las posiciones posibles. Esta práctica separa explícitamente los escenarios inicio, centro, final y ausente, de modo que las mediciones representan esos casos condicionados y no una distribución real de consultas de clientes (McGeoch, 2012; Cormen et al., 2022; Zanabria Gálvez, 2026).

\subtitulo{2.9. Búsqueda lineal}
La búsqueda lineal recorre secuencialmente una colección hasta hallar la clave o agotar los elementos. Su principal ventaja es que no requiere que los datos estén ordenados ni una preparación previa. En un arreglo, el procedimiento es directo y aprovecha un patrón de acceso secuencial. Su mejor caso requiere un sondeo, mientras que el peor caso requiere $n$ sondeos; bajo supuestos uniformes de éxito, su costo esperado sigue siendo lineal. La implementación iterativa utiliza $\Theta(1)$ memoria auxiliar (Cormen et al., 2022; Skiena, 2020; Black, 2020).

En esta práctica se modela un recorrido equivalente al código: inicializar $n$, índice $i=0$ y contador $ops=0$; mientras $i<n$, examinar $a[i]$, incrementar $ops$, comparar con la clave y avanzar si no coincide. Cada operación cuesta 1; incrementar una variable cuesta 2 (suma y asignación). Se consideran obtener/asignar la longitud almacenada y retornar el par como operaciones abstractas unitarias. Este convenio hace trazable la función de costo y, aunque otro convenio produciría constantes diferentes, el orden lineal no cambia (Cormen et al., 2022; Levitin, 2012; Zanabria Gálvez, 2026).

\begin{center}\small
\begin{tabular}{p{8.4cm}rr}
\toprule
Operación & Ausente & Éxito en visita $k$\\\midrule
Inicializaciones de $n$, $i$ y $ops$ & 3 & 3\\
Control $i<n$ & $n+1$ & $k$\\
Acceso al elemento & $n$ & $k$\\
Incremento de $ops$: suma y asignación & $2n$ & $2k$\\
Comparación de igualdad & $n$ & $k$\\
Incremento de $i$: suma y asignación & $2n$ & $2(k-1)$\\
Retorno & 1 & 1\\\midrule
Total & $7n+5$ & $7k+2$\\\bottomrule
\end{tabular}\end{center}

Así, $T_{\rm ausente}(n)=7n+5$ y, para $n\ge1$, $7n\le7n+5\le12n$, por lo que el peor caso es $\Theta(n)$. El contador almacenado en los CSV mide únicamente sondeos, no todas las operaciones del convenio RAM: vale $n$ en ausencia y $k$ cuando el éxito ocurre en la visita $k$. Esta separación permite comparar la predicción estructural con el tiempo real sin presentar ambas métricas como si fueran equivalentes (Cormen et al., 2022; McGeoch, 2012; Zanabria Gálvez, 2026).

\subtitulo{2.10. Búsqueda binaria}
La búsqueda binaria requiere una colección ordenada y compara la clave con el elemento central del intervalo activo. Si no coincide, descarta la mitad que no puede contener la clave. La secuencia aproximada de tamaños es $n,n/2,n/4,n/8,\ldots$; después de $k$ divisiones queda $n/2^k$. Al imponer $n/2^k\le1$ se obtiene $k\ge\log_2 n$, lo que explica el crecimiento logarítmico del número de iteraciones (Cormen et al., 2022; Sedgewick \& Wayne, 2011; Black, 2022).

Su mejor caso es $\Theta(1)$ cuando la primera posición central contiene la clave, mientras que el peor caso es $\Theta(\log n)$. Con un arreglo, la implementación iterativa utiliza $\Theta(1)$ memoria auxiliar. La ventaja asintótica frente a búsqueda lineal puede ser enorme para entradas grandes, pero depende de que los datos ya estén ordenados o de que el costo de ordenarlos pueda justificarse. Además, el acceso no secuencial y los costos concretos de la implementación influyen en el tiempo observado sin modificar la complejidad teórica (Cormen et al., 2022; Skiena, 2020; Levitin, 2012).

\subtitulo{2.11. Costo de preparación y amortización}
Comparar una búsqueda lineal sobre datos desordenados con una búsqueda binaria sobre datos ordenados exige contabilizar el costo de preparación. Si una colección debe ordenarse antes de la primera búsqueda, el costo total puede modelarse como $C_B(q)=S+qB$, donde $S$ es el costo de ordenar, $B$ el costo de una búsqueda binaria y $q$ la cantidad de consultas. Para lineal, si no existe preparación equivalente, puede modelarse $C_L(q)=qL$. La comparación relevante no es solamente $L$ frente a $B$, sino el costo acumulado de la estrategia completa (Cormen et al., 2022; Skiena, 2020; McGeoch, 2012).

Cuando $L>B$, la estrategia con ordenamiento empieza a ser favorable si $q>S/(L-B)$. Este umbral depende de los tiempos observados, del tipo de consulta y de la estabilidad de la colección; si los datos cambian continuamente, el costo de mantener el orden puede alterar la decisión. Por eso la búsqueda binaria no debe presentarse como universalmente superior: para una única consulta sobre datos desordenados, recorrer directamente puede ser más razonable; para miles de consultas sobre una colección estable, la preparación puede amortizarse (Levitin, 2012; Skiena, 2020; McGeoch, 2012).

\subtitulo{2.12. Complejidad teórica, rendimiento empírico y reproducibilidad}
La complejidad asintótica describe cómo escala una función de costo; el tiempo real incorpora además constantes, arquitectura, compilador, intérprete, representación de datos, caché, predicción de saltos, sistema operativo y carga concurrente. Por ello, un resultado experimental no ``demuestra'' una cota matemática. Lo correcto es verificar si las tendencias observadas son consistentes con la predicción teórica y analizar las discrepancias. Esta complementariedad entre análisis formal y experimentación es central en la evaluación rigurosa de algoritmos (McGeoch, 2012; Cormen et al., 2022; Skiena, 2020).

La reproducibilidad exige documentar datos, versiones, semillas, región cronometrada, número de repeticiones y estadísticos. En Python, \texttt{perf\_counter\_ns()} proporciona un contador de alta resolución apropiado para medir intervalos; en C++, \texttt{steady\_clock} es monotónico y está diseñado para medición de duraciones. Aun así, expresar un resultado en nanosegundos no implica precisión física de un nanosegundo, y las mediciones muy pequeñas pueden estar dominadas por sobrecarga del reloj o ruido. Por eso esta práctica conserva datos crudos, reporta mediana además de promedio y separa carga, ordenamiento y búsqueda (Python Software Foundation, s. f.; ISO/IEC JTC1/SC22/WG21, s. f.; McGeoch, 2012).

En conjunto, estos doce núcleos muestran que elegir un algoritmo no consiste únicamente en memorizar una notación. Es necesario definir la entrada, modelar operaciones, distinguir cotas de escenarios, considerar requisitos de los datos y comprobar empíricamente cómo se comporta una implementación real. La búsqueda lineal y la binaria son un caso de estudio especialmente útil porque presentan crecimientos teóricos diferentes y permiten observar cómo la preparación, el número de consultas y el entorno modifican la decisión práctica sin alterar las cotas matemáticas (Cormen et al., 2022; Levitin, 2012; McGeoch, 2012).
